% GENERATED FILE. Edit canonical lessons, then run npm run generate:books.
% canonical-source-hash: fnv1a64:489b291653bc3702
% canonical-lessons: JA-C161-sentaku, JA-C161-nureru, JA-C161-hirogeru, JA-C161-kawaku, JA-C161-kawakasu

\chapter{Washing, Spreading and Drying}
\label{ch:ja-washing-spreading-and-drying}
\hlchaptermodality{161}

\begin{chapteropening}
\textbf{By the end of this chapter:} \emph{I can talk about doing the washing and drying it.}

Complete the last lesson of chapter 161: I can talk about doing the washing and drying it.
\end{chapteropening}

\section[\texorpdfstring{\ja{せんたく} (sentaku)}{せんたく (sentaku)}]{\ja{せんたく} (sentaku) — washing clothes; the washing}

\label{lesson:JA-C161-sentaku}



\begin{quote}
Before the new one: say the Japanese for to make dirty, then the Japanese for dirty.
\end{quote}

\subsection*{You'll want to know: \ja{せんたく}}
\textbf{\ja{せんたく}} — \emph{sentaku} — \textquotedblleft{}washing clothes; the washing\textquotedblright{}.

\begin{quote}
\textbf{\ja{せんたくは} \ja{あにの} \ja{しごとです。}} — \emph{sentaku wa ani no shigoto desu} — \textquotedblleft{}The washing is my elder brother's job.\textquotedblright{}
\end{quote}

\begin{cousinweb}
\textbf{\ja{せんたく}} (\emph{sentaku}) is Sino-Japanese (kango): two Chinese characters read with their on'yomi, readings borrowed from Chinese. They are \emph{sen}, \textquotedblleft{}wash\textquotedblright{}, whose sound matches the rhyme book's \emph{senX} although it gives \textquotedblleft{}wash\textquotedblright{} under \emph{sejX}, and \emph{taku}, \textquotedblleft{}rinse\textquotedblright{} (Middle Chinese \emph{draewk}) (KANJIDIC2; Middle Chinese in Baxter's transcription).
\end{cousinweb}

\begin{grammarlens}[title={What you've built}]
One more word for everyday life at A2.
\end{grammarlens}

\subsection*{Your turn}
\begin{itemize}
\raggedright
  \item \emph{Say it:} \emph{sentaku}
  \item \emph{Say it:} \emph{sentaku}, once more
  \item \emph{Say it:} \emph{sentaku wa ani no shigoto desu} — The washing is my elder brother's job.
  \item \emph{Recall:} say the Japanese for to make dirty, then the Japanese for dirty, then say \emph{sentaku} again
\end{itemize}

\begin{tcolorbox}[breakable,colback=teal!4,colframe=teal!35!black,title={Before you move on}]
What does \ja{せんたく} mean? (\textquotedblleft{}Washing clothes; the washing\textquotedblright{}.) Say it once more.
\end{tcolorbox}
\section[\texorpdfstring{\ja{ぬれる} (nureru)}{ぬれる (nureru)}]{\ja{ぬれる} (nureru) — to get wet}

\label{lesson:JA-C161-nureru}



\begin{quote}
Before the new one: say the Japanese for dirty, then the Japanese for washing clothes.
\end{quote}

\subsection*{You'll want to know: \ja{ぬれる}}
\textbf{\ja{ぬれる}} — \emph{nureru} — \textquotedblleft{}to get wet\textquotedblright{}.

\begin{quote}
\textbf{\ja{あめで} \ja{くつが} \ja{ぬれます。}} — \emph{ame de kutsu ga nuremasu} — \textquotedblleft{}My shoes get wet in the rain.\textquotedblright{}
\end{quote}

Polite \textbf{\ja{ぬれます}}, \emph{nuremasu}; an ichidan verb.

\begin{cousinweb}
\textbf{\ja{ぬれる}} (\emph{nureru}) is a native Japanese word (wago), inherited rather than borrowed from Chinese; its kanji was attached later as a kun'yomi, and its older origin is not settled (JMdict, KANJIDIC2).
\end{cousinweb}

\begin{grammarlens}[title={What you've built}]
One more word for everyday life at A2.
\end{grammarlens}

\subsection*{Your turn}
\begin{itemize}
\raggedright
  \item \emph{Say it:} \emph{nureru}
  \item \emph{Say it:} \emph{nureru}, once more
  \item \emph{Say it:} \emph{ame de kutsu ga nuremasu} — My shoes get wet in the rain.
  \item \emph{Recall:} say the Japanese for dirty, then the Japanese for washing clothes, then say \emph{nureru} again
\end{itemize}

\begin{tcolorbox}[breakable,colback=teal!4,colframe=teal!35!black,title={Before you move on}]
What does \ja{ぬれる} mean? (\textquotedblleft{}To get wet\textquotedblright{}.) Say it once more.
\end{tcolorbox}
\section[\texorpdfstring{\ja{ひろげる} (hirogeru)}{ひろげる (hirogeru)}]{\ja{ひろげる} (hirogeru) — to spread out; to unfold}

\label{lesson:JA-C161-hirogeru}



\begin{quote}
Before the new one: say the Japanese for washing clothes, then the Japanese for to get wet.
\end{quote}

\subsection*{You'll want to know: \ja{ひろげる}}
\textbf{\ja{ひろげる}} — \emph{hirogeru} — \textquotedblleft{}to spread out; to unfold\textquotedblright{}.

\begin{quote}
\textbf{\ja{ふとんを} \ja{ゆかに} \ja{ひろげます。}} — \emph{futon o yuka ni hirogemasu} — \textquotedblleft{}I spread the futon out on the floor.\textquotedblright{}
\end{quote}

Polite \textbf{\ja{ひろげます}}, \emph{hirogemasu}; an ichidan verb.

\begin{cousinweb}
\textbf{\ja{ひろげる}} (\emph{hirogeru}) is a native Japanese word (wago), inherited rather than borrowed from Chinese; its kanji was attached later as a kun'yomi, and its older origin is not settled (JMdict, KANJIDIC2).
\end{cousinweb}

\begin{grammarlens}[title={What you've built}]
One more word for everyday life at A2.
\end{grammarlens}

\subsection*{Your turn}
\begin{itemize}
\raggedright
  \item \emph{Say it:} \emph{hirogeru}
  \item \emph{Say it:} \emph{hirogeru}, once more
  \item \emph{Say it:} \emph{futon o yuka ni hirogemasu} — I spread the futon out on the floor.
  \item \emph{Recall:} say the Japanese for washing clothes, then the Japanese for to get wet, then say \emph{hirogeru} again
\end{itemize}

\begin{tcolorbox}[breakable,colback=teal!4,colframe=teal!35!black,title={Before you move on}]
What does \ja{ひろげる} mean? (\textquotedblleft{}To spread out; to unfold\textquotedblright{}.) Say it once more.
\end{tcolorbox}
\section[\texorpdfstring{\ja{かわく} (kawaku)}{かわく (kawaku)}]{\ja{かわく} (kawaku) — to get dry}

\label{lesson:JA-C161-kawaku}



\begin{quote}
Before the new one: say the Japanese for to get wet, then the Japanese for to spread out.
\end{quote}

\subsection*{You'll want to know: \ja{かわく}}
\textbf{\ja{かわく}} — \emph{kawaku} — \textquotedblleft{}to get dry\textquotedblright{}.

\begin{quote}
\textbf{\ja{なつは} \ja{せんたくものが} \ja{よく} \ja{かわきます。}} — \emph{natsu wa sentakumono ga yoku kawakimasu} — \textquotedblleft{}In summer the laundry dries well.\textquotedblright{}
\end{quote}

Polite \textbf{\ja{かわきます}}, \emph{kawakimasu}; a godan verb.

\begin{cousinweb}
\textbf{\ja{かわく}} (\emph{kawaku}) is a native Japanese word (wago), inherited rather than borrowed from Chinese; its kanji was attached later as a kun'yomi, and its older origin is not settled (JMdict, KANJIDIC2).
\end{cousinweb}

\begin{grammarlens}[title={What you've built}]
One more word for everyday life at A2.
\end{grammarlens}

\subsection*{Your turn}
\begin{itemize}
\raggedright
  \item \emph{Say it:} \emph{kawaku}
  \item \emph{Say it:} \emph{kawaku}, once more
  \item \emph{Say it:} \emph{natsu wa sentakumono ga yoku kawakimasu} — In summer the laundry dries well.
  \item \emph{Recall:} say the Japanese for to get wet, then the Japanese for to spread out, then say \emph{kawaku} again
\end{itemize}

\begin{tcolorbox}[breakable,colback=teal!4,colframe=teal!35!black,title={Before you move on}]
What does \ja{かわく} mean? (\textquotedblleft{}To get dry\textquotedblright{}.) Say it once more.
\end{tcolorbox}
\section[\texorpdfstring{\ja{かわかす} (kawakasu)}{かわかす (kawakasu)}]{\ja{かわかす} (kawakasu) — to dry (something)}

\label{lesson:JA-C161-kawakasu}



\begin{quote}
Before the new one: say the Japanese for to spread out, then the Japanese for to get dry.
\end{quote}

\subsection*{You'll want to know: \ja{かわかす}}
\textbf{\ja{かわかす}} — \emph{kawakasu} — \textquotedblleft{}to dry (something)\textquotedblright{}.

\begin{quote}
\textbf{\ja{ふとんを} \ja{かわかします。}} — \emph{futon o kawakashimasu} — \textquotedblleft{}I dry the futon.\textquotedblright{}
\end{quote}

Polite \textbf{\ja{かわかします}}, \emph{kawakashimasu}; a godan verb. Its intransitive partner is \textbf{\ja{かわく}}, \emph{kawaku}.

\begin{cousinweb}
\textbf{\ja{かわかす}} (\emph{kawakasu}) is a native Japanese word (wago), inherited rather than borrowed from Chinese, and built on the same root as \emph{kawaku}, \textquotedblleft{}to get dry\textquotedblright{}; its kanji was attached later as a kun'yomi (JMdict, KANJIDIC2).
\end{cousinweb}

\begin{grammarlens}[title={What you've built}]
One more word for everyday life at A2.
\end{grammarlens}

\subsection*{Your turn}
\begin{itemize}
\raggedright
  \item \emph{Say it:} \emph{kawakasu}
  \item \emph{Say it:} \emph{kawakasu}, once more
  \item \emph{Say it:} \emph{futon o kawakashimasu} — I dry the futon.
  \item \emph{Recall:} say the Japanese for to spread out, then the Japanese for to get dry, then say \emph{kawakasu} again
\end{itemize}

\begin{tcolorbox}[breakable,colback=teal!4,colframe=teal!35!black,title={Before you move on}]
What does \ja{かわかす} mean? (\textquotedblleft{}To dry (something)\textquotedblright{}.) Say it once more.
\end{tcolorbox}
